\scanpage{109}
\deleted{On va \deleted{\unsure{désormais|d'ailleurs}} supposer que $J$ a un plus petit élément $e_J$, et $d(J)=d$ sa codim.}

Alors la multiplicité $D_J$ est lisse, \deleted{de} et le morphisme canonique
\[
(153)\qquad D_J\longrightarrow X
\]
est une imm. \unsure{locale|local} de codim $d$.
Elle se factorise par $D$, sauf dans le cas où $d(I)$\inserted{$=0$ i.e. $I$} est réduit à un seul élément, de codim. nulle, auquel cas $D_J=\deleted{D_0}=X$.

Si $J$ est formé d'un seul élément, de codim $d$, on retrouve $D_d$; S'il est formé de deux éléments, \deleted{le plus} $e_J$ de codim $d$ et $e'$ de codim $1$, on retrouve $\widetilde D_d$.

Si $J$ est de dim totale $d=d(J)$, \unsure{on peut} \deleted{\unsure{retrouver}} reconstituer $D_J$ à \unsure{l'aide de} $D_d$ et de $\widetilde D_d$, via le \unsure{revêtement} principal (torseur sous $\unsure{\mathfrak S|\mathcal S}_d$) \unsure{sur} $D_d$ associé à $\widetilde D_d$, soit $\overset{\deleted{\unsure{\sim}}}{D'_d}$, qui s'identifie aussi à $D_{\mathbf I(d)}$, où $\mathbf I(d)$ l'ens ordonné
\diagram{\begin{tikzcd}[ampersand replacement=\&, column sep=small]
\underset{d}{\bullet}\arrow[r,dash]\&\underset{d-1}{\bullet}\arrow[r,dash]\&\underset{\unsure{2}}{\bullet}\arrow[r,dash]\&\underset{1}{\bullet}\arrow[r,dash]\&\underset{0}{\bullet}
\end{tikzcd}}
(degré maximal de codim $d$). En effet, considérons
\margin{Mais l'expression de $\overset{\deleted{\unsure{\sim}}}{D'_d}$ comme $D_{I_d}$ \unsure{ne met pas en évidence les opérations de|ne vaut pas une évidence, car opération de} $\unsure{\mathfrak S|\mathcal S}_d$ \unsure{sur} $\widetilde D_d$ — \unsure{pour ceci}; il faut \unsure{prendre} \illegible{2} l'ens ordonné $\unsure{K}_d$, \unsure{ayant} un \deleted{\unsure{élément}} plus petit élément de codim $d$, et $d$ éléments de codim $1$\ldots}
\margin{les \unsure{entiers} sont les \unsure{sommets identifiés par} leur codim.}
\margin{108}
